\label{sec:conclusion}

M.4 operationalizes commuting shed community membership as an edge weight
signal through the Census Bureau's LODES Origin-Destination file.
The weighted Jaccard similarity of commuting destination distributions
measures the overlap between two tracts' daily economic spheres: tracts
whose workers commute to the same employment centers share a functional
economic area and receive heavy edges that METIS avoids cutting.

The signal fills a gap that M.1 (economic character) leaves open.
M.1 identifies employment centers---where jobs are concentrated.
M.4 identifies which residential communities belong to which employment
center's catchment area.
In polycentric metropolitan areas, where multiple large employment centers
compete for workers from the same residential zones, M.4 provides the
signal that correctly assigns residential tracts to their primary functional
economic community.

\paragraph{Relationship to M.1.}
The complementarity between M.1 and M.4 is a design feature of the M-track
series.
M.1 fires at residential-employment center boundaries (the RTP campus
tract has high JPR, its residential neighbors have low JPR: a strong
M.1 cut seam).
M.4 fires at commuting shed boundaries (the RTP residential catchment
area has high Triangle destination overlap; Charlotte residential tracts
have low Triangle destination overlap: a strong M.4 cut seam further from
the employment center).
Together, M.1 and M.4 define the economic community at two different
geographic scales: the employment cluster boundary (M.1) and the
residential catchment boundary (M.4).
The M.8 composite index \citep{m8composite} assigns M.1 a weight of 0.20
and M.4 a weight of 0.10, reflecting the judgment that M.4 is a meaningful
but more indirect community signal than the direct economic character
differentiation of M.1.

\paragraph{Commuting data and the daily life argument.}
The commuting shed argument for communities of interest is grounded in
the observation that daily economic behavior---not formal administrative
membership or physical proximity---defines the economic community most
people live in.
The farmer who drives an hour to a market town, the suburban worker who
commutes to an office park, the airport employee who drives the same
route to the same terminal every morning: these are the daily economic
communities that redistricting commissions are asked to preserve.
The LODES OD data makes these communities measurable for the first time
in an algorithmic redistricting system.

\paragraph{Phase 2 commitments.}
Phase~2 will complete the following quantitative results, all currently
marked ${}^{\dagger}$:
\begin{enumerate}
  \item Milwaukee suburban co-assignment rate (SCR) for WI-8,
    comparing commuting-shed weights to geographic baseline.
    Prediction: $\mathrm{SCR}_{\mathrm{commute}} \geq 0.70$ vs.\
    $\mathrm{SCR}_{\mathrm{geographic}} \leq 0.50$.
  \item Research Triangle / Charlotte cross-metro co-assignment rate (CMR)
    for NC-14.
    Prediction: $\mathrm{CMR}_{\mathrm{commute}} < 0.05$ vs.\
    $\mathrm{CMR}_{\mathrm{geographic}} \leq 0.15$.
  \item Weighted Jaccard wall time for NC.
    Prediction: $< 60$ seconds on 8 workers.
  \item Proportionality gap comparison for WI-8 and NC-14 (secondary metric;
    predicted to show improvement for NC and null for WI).
\end{enumerate}

Phase~2 depends on LODES OD data download via
\texttt{bisect fetch --type lodes-od --year 2020}.
The same LEHD infrastructure used for M.1 WAC data handles this download;
the OD file uses the same authentication and the same Parquet caching
pipeline.
